\documentclass{tikzposter} %Options for format can be included here
\geometry{paperwidth=850mm, paperheight=1600mm}
\makeatletter
\setlength{\TP@visibletextwidth}{\textwidth-2\TP@innermargin}
\setlength{\TP@visibletextheight}{\textheight-2\TP@innermargin}
\makeatother
\usepackage{amsmath}
% \usepackage{eucal}
\usepackage{bm}
\usepackage{mathrsfs}
\usepackage{amssymb}
\usepackage{dsfont}
\usepackage{enumitem}
\usepackage{parskip}
\usepackage{amsfonts}
\usepackage{verbatim}
\usepackage{nicefrac}
\usepackage{xcolor}
\usepackage{mathtools}
\newcommand{\mc}{\mathcal}
\renewcommand{\P}{\mathbb{P}}
\newcommand{\E}{\mathbb{E}}
\newcommand\Eb[1]{\mathbb{E}\big[#1\big]}
\DeclarePairedDelimiter{\floor}{\lfloor}{\rfloor}
\DeclarePairedDelimiter{\ceil}{\lceil}{\rceil}
\DeclareMathOperator{\Int}{int}
\DeclareMathOperator{\cl}{cl}
\DeclareMathOperator{\Var}{Var}
\DeclareMathOperator{\degree}{deg}
\newcommand\restr[2]{{% we make the whole thing an ordinary symbol
  \left.\kern-\nulldelimiterspace % automatically resize the bar with \right
  #1 % the function
  \vphantom{\big|} % pretend it's a little taller at normal size
  \right|_{#2} % this is the delimiter
  }}
\newcommand\leftopen[2]{\ensuremath{(#1,#2]}}
\newcommand\rightopen[2]{\ensuremath{[#1,#2)}}


\newtheorem{theorem}{Theorem}
\newtheorem{lemma}[theorem]{Lemma}
\newtheorem{corollary}{Corollary}

\newtheorem{definition}{Definition}

\newtheorem{remark}{Remark}
\newtheorem{claim}{Claim}
\newtheorem{case}{Case}

\definecolor{nbYellow}{HTML}{FCF434}
\definecolor{nbPurple}{HTML}{9C59D1}
\definecolor{nbBlack}{HTML}{2C2C2C}
\definecolor{tBlue}{HTML}{5BCEFA}
\definecolor{tPink}{HTML}{F5A9B8}
\definecolor{bp1}{HTML}{D60270}
\definecolor{bp2}{HTML}{9B4F96}
\definecolor{bp3}{HTML}{0038A8}
\definecolor{pcs1}{HTML}{B300B3}
\definecolor{pcs2}{HTML}{54007D}
\definecolor{pcs3}{HTML}{B30086}
\definecolor{pcs4}{HTML}{3C00B3}
\definecolor{pcs5}{HTML}{2A007D}

\definecolorstyle{NewColour} {
  \definecolor{c1}{named}{nbBlack}
  \definecolor{c2}{named}{nbPurple}
  \definecolor{c3}{named}{nbYellow}
}{
  % Background Colors
  \colorlet{backgroundcolor}{black!10}
  \colorlet{framecolor}{black}
  % Title Colors
  \colorlet{titlefgcolor}{black}
  \colorlet{titlebgcolor}{black!10}
  % Block Colors
  \colorlet{blocktitlebgcolor}{c1}
  \colorlet{blocktitlefgcolor}{white}
  \colorlet{blockbodybgcolor}{white}
  \colorlet{blockbodyfgcolor}{black}
  % Innerblock Colors
  \colorlet{innerblocktitlebgcolor}{c2!80}
  \colorlet{innerblocktitlefgcolor}{black}
  \colorlet{innerblockbodybgcolor}{c2!50}
  \colorlet{innerblockbodyfgcolor}{black}
  % Note colors
  \colorlet{notefgcolor}{black}
  \colorlet{notebgcolor}{c3!50}
  \colorlet{notefrcolor}{c3!70}
}

\defineblockstyle{NewBlock}{
  titlewidthscale=1, bodywidthscale=1, titleleft,
  titleoffsetx=0pt, titleoffsety=0pt, bodyoffsetx=0pt, bodyoffsety=0pt,
  bodyverticalshift=0pt, roundedcorners=0, linewidth=0pt, titleinnersep=1cm,
  bodyinnersep=1cm
}{
  \ifBlockHasTitle%
  \draw[draw=none, fill=blocktitlebgcolor]
  (blocktitle.south west) rectangle (blocktitle.north east);
  \fi%
  \draw[draw=none, fill=blockbodybgcolor] %
  (blockbody.north west) [rounded corners=30] -- (blockbody.south west) --
  (blockbody.south east) [rounded corners=0]-- (blockbody.north east) -- cycle;
}

% Choose Layout
\usecolorstyle{NewColour}
\usebackgroundstyle{Default}
\usetitlestyle{Filled}
\useblockstyle{NewBlock}
\useinnerblockstyle[roundedcorners=0.2]{Default}
\usenotestyle[roundedcorners=0]{Default}

\settitle{\centering \color{titlefgcolor} {\Large \@title \, -- \, \@author}}

% Title, Author, Institute
\title{Optimal Control}
\author{Ike Glassbrook}
\begin{document}
\maketitle
\begin{columns}
  \column{0.5}
  \block{Discrete-time control and MDPs}{
    \begin{definition}
    \ A discrete-time control problem is formed of the following:
    \begin{itemize}[label=$\cdot$]
            \item \ A control process $(U_t)_{t \in [T]}$ in some $\mc{U}$.
            \item \ A state process $(X^U_t)_{t \in [T]}$ in $\mc{X} \subseteq \mathbb{R}^d$, a function of prior controls.
            \item \ An instantaneous cost function $g : [T] \times \mc{X} \times \mc{U} \to \mathbb{R}$.
    \end{itemize}
    An \emph{optimal control} is a $U^{\ast} = (U_t^{\ast})$ which minimises
    \begin{align*}
      J(x, U) = \sum_{t \in [T]} g(t, X_t^U, U_t)
    \end{align*}
    where $X_0 = x$.
    \end{definition}
    \hphantom{}

    To solve this, we decompose the problem into minimising each of the `remaining-cost' problems. With
    \begin{align*}
      J(t, x, U) = \sum_{s \ge t : s \in [T]} g(s, X_s^{t,x,U}, U_s)
    \end{align*}
    the remaining cost from time $t$ onwards given $X_t = x$, we note that
    \begin{align*}
      J(t, x, U) &= g(t, x, U_t) + J(t+1, X^{t, x, U}_{t+1}, U)
    \end{align*}
    and with $V(t, x) := \inf_U J(t, x, U)$, it is clear by the independence of $J$ from prior decisions and states, and the assumed Markov property of $(X_t)$, $X_{t+1} = f(t, X_t, U_t)$ that
    \begin{align*}
      V(t, x) = \inf_{U_t} \left\{g(t, x, U_t) + V\big(t+1, f(t, x, U_t)\big)\right\}
    \end{align*}
    which is the Bellman equation.






  }
  \block{Numerical methods and Reinforcement Learning}{

  }
  \block{Continuous Deterministic Control}{

  }
  \block{Continuous Stochastic Control}{

  }

\end{columns}
\end{document}
